% STATUS: part written 2026-11-04 (Miers 1978 read in full); the n-Lie half
% added 2026-11-06 with one source unobtainable. Abdullaev 1992 (Uzbek. Mat.
% Zh.) has no open copy and no shadow-library copy; the 2026-11-06 searches
% through papergraph found only later work on n-Lie derivations, which is a
% different literature and is not substituted here. The notion is named and
% the source cited, with no statement attributed to it.
\section{Lie triple and Lie n-derivations}\label{sec:lie-type:triple-and-n}

Sections~\ref{sec:lie-type:martindale} and~\ref{sec:lie-type:bresar} treated the
commutator of two elements. The next kind iterates the bracket one step further.
Miers defined it for a von Neumann algebra and proved that the same kind of
decomposition as for Lie derivations holds \cite[p.~57]{Miers1978}.

\begin{definition}[{\cite[p.~57]{Miers1978}}]\label{def:lie:miers}
Let $M$ be a von Neumann algebra. An
associative derivation of $M$ is an additive map $d\colon M\to M$ satisfying
$d(xy)=d(x)y+xd(y)$. A Lie derivation of $M$ is an additive map $L\colon M\to M$
satisfying $L([x,y])=[L(x),y]+[x,L(y)]$. Generalizing both of these definitions,
a \emph{Lie triple derivation} of $M$ is an additive map $T\colon M\to M$ with the
property that
\[
T[[x,y],z]=[[T(x),y],z]+[[x,T(y)],z]+[[x,y],T(z)]\qquad\text{for all }x,y,z\in M.
\]
\end{definition}

Miers proves his main result together with the Lie derivation form of
Theorem~\ref{thm:lie:bresar-standard} for von Neumann algebras.

\begin{theorem}[{\cite[abstract and p.~58]{Miers1978}}]\label{thm:lie:miers-lielet}
If $L\colon M\to M$ is a Lie derivation of the von Neumann algebra $M$, then
there exists an associative derivation $d\colon M\to M$ and an additive map
$t\colon M\to \mathrm{center}(M)$ such that $d(x)+t(x)=L(x)$ for all $x\in M$.
If $T\colon M\to M$ is a Lie triple derivation of the von Neumann algebra $M$,
then there exists an associative derivation $d\colon M\to M$ and a linear map
$\lambda\colon M\to \mathrm{center}(M)$ with $\lambda[[x,y],z]=0$ for all
$x,y,z\in M$ such that $d(x)+\lambda(x)=T(x)$ for all $x\in M$.
\end{theorem}

Miers adds a form that respects the involution.

\begin{theorem}[{\cite[Theorem~4, p.~60]{Miers1978}}]\label{thm:lie:miers-star}
If $T$ is a Lie $*$-triple derivation of $M$, so that $T(x^{*})=T(x)^{*}$, then
there exists an associative $*$-derivation $d$ and a linear map
$\lambda\colon M\to Z(M)$ with $\lambda[[x,y],z]=0$ for all $x,y,z\in M$ and
$\lambda(x^{*})=-\lambda(x)^{*}$ such that $T(x)=d(x)+\lambda(x)$.
\end{theorem}

The next kind after the triple derivation maps into $n$-fold commutators:
an \emph{$n$-Lie derivation} is a map of the algebra, taken in the operator
setting by Abdullaev \cite{Abdullaev1992}.
His paper states the theorem for von Neumann algebras.
No copy of it has been obtained, so this book names the notion and cites the
source without quoting a statement from it.
The later literature on $n$-Lie derivations of matrix and triangular algebras
is recorded in Chapter~\ref{ch:other-systems}.

On a Lie triple system the Jordan condition is enough. Najati shows under some
conditions that every Jordan triple $(\phi,\psi)$-derivation of a Lie triple
system is a $(\phi,\psi)$-derivation, and in particular that every Jordan triple
$\theta$-derivation is a $\theta$-derivation
\cite[Abstract, p.~143]{Najati2009LieTripleSystems}.
